\section*{Chapter \ref{ch:iv1:timed-test-7} --- Timed Test 7}
\markboth{\protect\omTocMark{}{Answers: Timed Test 7}}{\protect\omTocMark{}{Answers: Timed Test 7}}

\begin{solution}{iq:iv1:timed-test-7:1}
The first is 9 rolls, and the second is 11.

The wait splits into stages\index{coupon collecting}. At first, either of the 2 wanted faces will do, so the wait is $\frac62 = 3$. Then one face is left, and the wait is 6. The sum is 9.

For three faces, the stages are $\frac63 + \frac62 + \frac61 = 2 + 3 + 6 = 11$.

An exact recursion and a simulation agree. The last face takes more than half of the wait here. It takes the largest single share in any collecting problem, though less than half once you want four faces or more.

A second route for the first figure uses inclusion--exclusion on the maximum of two waits: $6 + 6 - 3 = 9$, where 3 is the wait for the first of the two.

\iqwrong{Adding the separate waits, $6 + 6 = 12$, which double-counts the rolls during which both faces are still missing.}
\end{solution}

\begin{solution}{iq:iv1:timed-test-7:2}
The probability is about $0.554$.

Expand $(q + p)^n$ and $(q - p)^n$\index{parity trick}. Adding them cancels the odd powers of $p$, so
\[
P(\text{even}) = \frac{1 + (1 - 2p)^n}{2} = \frac{1 + 0.8^{10}}{2} = 0.5537.
\]

Limits: at $p = \frac12$ the formula gives exactly $\frac12$, and at $p = 0$ it gives 1, since 0 heads is even. As $n$ grows, it tends to $\frac12$ for any $p$ strictly between 0 and 1. An exact binomial sum confirms the value.

The same result follows from a recursion: a parity bias of $b_n = P(\text{even}) - P(\text{odd})$ satisfies $b_{n + 1} = (1 - 2p)b_n$, with $b_0 = 1$.

\iqwrong{Summing the even binomial terms by hand and stopping after a few, which is slow and stops short of the value: the terms for 0 and 2 heads give only $0.54$.}
\end{solution}

\begin{solution}{iq:iv1:timed-test-7:3}
The variance is 7, and the slope is $\frac12$.

Expanding\index{variance of a sum},
\[
\operatorname{Var}(X + 2Y) = 1 + 4 + 2 \times 2 \times \frac12 = 7.
\]

The slope is $\operatorname{Cov}(Y, X + Y)/\operatorname{Var}(X + Y)$. The covariance is $\frac12 + 1 = \frac32$, and the variance is $1 + 1 + 1 = 3$. So the slope is $\frac12$.

There is a second route by symmetry. The slopes of $X$ and of $Y$ on $X + Y$ add up to 1, since $X + Y$ regressed on itself has slope 1. And $X$ and $Y$ have the same variance, so the two slopes are equal. That gives $\frac12$ for any correlation. A simulation confirms it.

\iqwrong{Leaving a factor out of the cross term. It is $2\operatorname{Cov}(X, 2Y) = 2 \times 2 \times \frac12 = 2$, and an answer of 6 means that a factor of 2 was dropped.}
\end{solution}

\begin{solution}{iq:iv1:timed-test-7:4}
The expected value is $\frac{35}{36} \approx 0.97$.

By symmetry, $\mathbb{E}\max(X - Y, 0) = \frac12\mathbb{E}|X - Y|$\index{symmetry}. The differences of 1, 2, 3, 4 and 5 have 10, 8, 6, 4 and 2 of the 36 outcomes. So
\[
\mathbb{E}|X - Y| = \frac{10 + 16 + 18 + 16 + 10}{36} = \frac{70}{36}.
\]
Half of this is $\frac{35}{36}$. An enumeration of the 36 outcomes confirms it.

This is the value of a very small option: you receive the excess of one die over the other, and at the money it has a positive value even though $\mathbb{E}(X - Y) = 0$.

\iqwrong{Answering 0, because $X - Y$ has mean 0. The maximum drops the negative values, and so a mean of 0 gives no such shortcut.}

\iqwhy{It is the plain form of the point behind every option: convexity turns variance into value.}
\end{solution}

\begin{solution}{iq:iv1:timed-test-7:5}
The first is $\frac{1}{11}$, and the second is $\frac16$.

Out of the 36 outcomes, 11 have at least one 6, and only one has two\index{conditional probability}. That gives $\frac{1}{11}$.

If the first die is known, the other is independent, and the chance is $\frac16$.

The difference lies in what is known: ``at least one'' keeps 11 equally likely outcomes, and ``this one'' keeps 6. An enumeration gives the same.

Which reading applies depends on how the information came. If you asked someone ``is there a 6?'', the first applies. If you saw one die, the second does. In practice, that step decides the answer, so it is worth asking aloud.

\iqwrong{Saying $\frac16$ for both, on the grounds that the other die is independent. In the first case, there is no ``other die'' to point to.}
\end{solution}

\begin{solution}{iq:iv1:timed-test-7:6}
The probability is $\frac38$.

The determinant $ad - bc$ is zero when $ad = bc$\index{determinant}. Each product is 1 in one case out of 4.

Both products are 0 in $3 \times 3 = 9$ of the 16 matrices, and both are 1 in 1 more. So 10 matrices are singular, and 6 are invertible, which gives $\frac{6}{16}$.

An enumeration of all 16 gives the same. Over the reals, the 6 are those with exactly one of the two products equal to 1. As the size grows, the chance of singularity for random binary matrices goes to 0, so this low chance is a small-size effect.

\iqwrong{Assuming that random matrices are almost always invertible. That holds for continuous entries. With binary entries, the coincidence $ad = bc$ is common.}
\end{solution}

\begin{solution}{iq:iv1:timed-test-7:7}
The probability is $\frac{99}{116} \approx 0.853$.

The likelihoods of two heads are $\frac14$, 1 and $\frac{9}{16}$\index{Bayes' rule}. Multiplied by 16, they are 4, 16 and 9, for a total of 29. So the posterior weights are $\frac{4}{29}$, $\frac{16}{29}$ and $\frac{9}{29}$.

The next flip is a head with chance
\[
\frac{4 \cdot \frac12 + 16 \cdot 1 + 9 \cdot \frac34}{29} = \frac{24.75}{29} = \frac{99}{116}.
\]
A simulation agrees. Before any flips, the chance was $\frac{2.25}{3} = \frac34$. Two heads raise it to $0.853$, a sensible move because evidence favours the two-headed coin but cannot rule the others out.

\iqwrong{Taking the most likely coin and answering 1. The posterior is a weighting of all three coins.}

\iqfollow{What if the two flips were H, T?}{The two-headed coin is out. The weights are $\frac14$ and $\frac{3}{16}$, so 4 and 3, and the next head has chance $\frac{4 \cdot \frac12 + 3 \cdot \frac34}{7} = \frac{17}{28}$.}
\end{solution}

\enlargethispage{-4\baselineskip}%
\begin{solution}{iq:iv1:timed-test-7:8}
The expected age is $\frac23$, and the expected total life is $\frac43$.

A random moment is more likely to fall in a long life. A bulb with life $\ell$ is seen with weight proportional to $\ell$\index{inspection paradox}. So the bulb in use has an expected life of
\[
\frac{\mathbb{E}L^2}{\mathbb{E}L} = \frac{4/3}{1} = \frac43,
\]
which is above the plain mean of 1.

Within that bulb's life, the moment is uniform, so the expected age is half the life: $\frac23$. And the expected remaining life is also $\frac23$. A simulation of a long series of bulbs agrees.

\iqwrong{Answering 1 and $\frac12$. These use a typical bulb, but sampling at a random moment picks the longer ones.}

\iqwhy{The same effect makes a passenger wait longer than half of the average bus gap.}
\end{solution}

\begin{solution}{iq:iv1:timed-test-7:9}
$\mathbb{E}[Z \mid Z > 1] \approx 1.52$, and the expected shortfall is about $2.67$.

For a standard normal, $\varphi^\prime(z) = -z\varphi(z)$, which gives $\int_a^\infty z\varphi(z)\,dz = \varphi(a)$. So\index{truncated normal}
\[
\mathbb{E}[Z \mid Z > a] = \frac{\varphi(a)}{P(Z > a)}.
\]
At $a = 1$, this is $\frac{0.242}{0.159} = 1.52$.

The expected shortfall is the same formula at the quantile: $\frac{0.0267}{0.01} = 2.67$. It is only about $15\%$ above the value at risk here, as the tail thins fast. A fatter tail gives a larger gap. A numerical integral confirms both.

A check at $a = 0$: the formula gives $\frac{0.399}{0.5} = 0.80$, which is the mean of $|Z|$, as it should be.

\iqwrong{Dividing by the wrong tail, or quoting the quantile itself. The answer must be above $a$, and that check catches the first error.}
\end{solution}

\begin{solution}{iq:iv1:timed-test-7:10}
The mean is 3, and the variance is $\frac{13}{6}$.

Each of the $\binom42 = 6$ pairs is out of order with chance $\frac12$. So the mean is 3\index{indicator variables}.

For the variance, build the permutation by inserting elements in order. Element $k$ lands in one of $k$ positions, and adds a uniform number on $\{0, \dots, k - 1\}$ of inversions, independent of the earlier ones. A uniform on $\{0, \dots, k - 1\}$ has variance $\frac{k^2 - 1}{12}$. So
\[
\operatorname{Var} = \frac{0 + 3 + 8 + 15}{12} = \frac{26}{12} = \frac{13}{6}.
\]
An enumeration of the 24 permutations confirms it. The general result is $n(n - 1)(2n + 5)/72$. The insertion step removes all the covariances from the indicator route.

\iqwrong{Adding the pair variances, $6 \times \frac14 = \frac32$. The pairs that share a position are not independent.}
\end{solution}

\begin{solution}{iq:iv1:timed-test-7:11}
The expected value is $\sqrt{\frac{0.5}{\pi}} \approx 0.399$.

Write $\max(X, Y) = \frac{X + Y}{2} + \frac{|X - Y|}{2}$\index{maximum as absolute value}. The first term has mean 0. The difference $X - Y$ is normal with variance $2 - 2\rho = 1$.

For a normal with variance $s^2$, $\mathbb{E}|\cdot| = s\sqrt{2/\pi}$. So the mean of the maximum is $\frac12\sqrt{2/\pi}\sqrt{2 - 2\rho} = \sqrt{(1 - \rho)/\pi}$. At $\rho = \frac12$, this is $0.399$.

Limit checks: at $\rho = 1$, the two are equal and the maximum has mean 0. At $\rho = 0$, it is $\frac{1}{\sqrt\pi} = 0.564$, the known value for two independent normals. A simulation returns $0.399$.

\iqwrong{Scaling the independent answer by $1 - \rho$, which gives $0.282$. The dependence runs through $\sqrt{1 - \rho}$.}
\end{solution}

\begin{solution}{iq:iv1:timed-test-7:12}
The probability is about $0.04$.

The sum has mean 350 and variance $100 \cdot \frac{35}{12} = 291.7$, so its standard deviation is about $17.1$\index{central limit theorem}.

With the continuity correction, the threshold is $379.5$, and $z = \frac{29.5}{17.1} = 1.73$. So $P \approx 0.042$. The sum of 100 dice is close enough to normal at this level, and an exact convolution gives the same to within $0.001$.

Without the correction, $z = 1.75$, and the tail is $0.040$. For a quick answer, that is close enough, though the correction costs one extra step for the extra digit.

\iqwrong{Using a variance of 25, as for 100 coin flips. A die's variance is near 3, which is a benchmark worth knowing.}
\end{solution}
